Cieľom technológií umelej inteligencie (\acrshort{ui}) je pomocou technických prostriedkov napodobniť biologické princípy ako ľudského organizmu, iných živočíchov a~rastlín, ale aj prírodných zákonitostí ako takých. Výskum v~oblasti \acrshort{ui} začal už v~prvej polovici 20.\ storočia, dlhé desaťročia boli však viaceré koncepty \acrshort{ui} len na úrovni matematickej teórie, keďže vtedy známe technické prostriedky nedisponovali takým výpočtovým výkonom, ktorý by postačoval na efektívnu implementáciu týchto návrhov. Preto najmä v~posledných desaťročiach, keď sa hranice výpočtových možností posúvajú závratným tempom ďalej, sme svedkami masívneho rozmachu v~oblasti výskumu a~aplikácií metód umelej inteligencie. Uplatnenie nachádzajú najmä v~riešení komplexných úloh, ktoré sú problémové na analytický opis a~riešenie konvenčnými algebraickými metódami.

Problematika fúzie dátových tokov z~viacerých snímačov rovnakého alebo rôzneho typu je tiež jednou z~možných aplikácií \acrshort{ui}. Cieľom je aj v~tomto prípade zabezpečiť stabilitu výsledného dátového toku, znížiť šum meraní či zamedziť negatívnym vplyvom skokovitej chyby meraní niektorého zo senzorov. Oproti klasickým metódam ako sú Kalmanove či časticové filtre, potenciál nástrojov \acrshort{ui} je v~rozpoznávaní zložitejších vzorov v~senzorických dátach, schopnosti spracúvať veľké množstvo zdrojov dát, ale tiež v~možnosti uskutočniť komplexnejší rozhodovací proces. Na to sa využívajú viaceré prístupy \acrshort{ui}, ako sú umelé neurónové siete a~ich modifikácie (hlboké, konvolučné, rekurzívne), strojové učenie a~jeho modifikácie (učenie s~posilňovaním, hlboké učenie) či fuzzy logika. V~nasledujúcich častiach si predstavíme základný princíp fungovania týchto algoritmov a~ukážeme aj príklady ich použitia v~konkrétnych aplikáciách na fúziu údajov pre lokalizáciu.